% res_catch_offset.tex -- DRAFT 23 Oct 2026 (thesis result U): the magnet caught off its slot, twin only.
% Sources: card docs/experiments/2026-10-23_catch_offset.md (critic 23 Oct answered; section 5 FROZEN 07:56 on Wesley's
% four answers, decisions.md DECIDED 2026-10-23 (7): the separation reported apart from the arm comparison (d_sep,
% d_task); t116 a grid cell x2; past ~2.5 cm the claim is worded as a model-error tolerance; the mechanism bar kept);
% its Results (12:41, lab batch 16, R2629-R2684, tools/catch_offset.py as frozen, last edited 00:52); decisions.md
% DECIDED 2026-10-23 "U after batch 16" (Wesley ~14:01): (1) the 6 owed runs fly on today's tree beside batch 16's 54
% counted runs, the table noting both trees; U re-flies whole on the final tree at the results freeze, like R; R2634 is
% a runner race (mocap staleness after /fleet/landed), VOID and re-flown; (2) along the rim the mechanism is NOT SHOWN
% (the scorer's pre-written code), radially in it is shown; the thesis says the slot arm folds in more gently along the
% rim, not that the model point does not matter.
% Runs (repeats a / b). Shared d = 0: measured R2629/R2657, slot R2630/R2658. Along the rim, measured: t02 R2631/R2659,
% t04 R2633/R2661, t06 R2635/R2663, t08 R2637 (VOID: weld 56.3 s after MAGNET ON, 14.3 s before the path)/R2665, t116
% R2639/R2667, t12 R2641/R2669 (R2669 never settled), t16 R2643/R2671 (separation FAIL 0.22/0.23 m). Along the rim, slot:
% t02 R2632/R2660, t04 R2634 (VOID, runner race after LANDED; scorer's code counts it FAIL)/R2662, t06 R2636 (FAIL 1, no
% weld, WELD_TIMEOUT 90 s)/R2664, t08 R2638/R2666, t116 R2640/R2668, t12 R2642/R2670, t16 R2644/R2672 (FAIL 1: weld 68.9 s
% after MAGNET ON; not re-flown). Radially in, measured: ri02 R2645/R2673, ri04 R2647/R2675, ri06 R2649/R2677, ri08
% R2651/R2679, ri12 R2653/R2681, ri16 R2655/R2683; slot: R2646/R2674, R2648/R2676, R2650/R2678, R2652/R2680, R2654/R2682,
% R2656/R2684. Tree be34f95 + patch cf1ab9ad, lab PC, 4 slots, RTF 0.20-0.25, world three_rigid_ground_rig_attach_L100,
% SIM_MOCAP_VEL_FILTER ema. The E-runs R2459-R2463 and the neighbours (u_meas_ta08, u_meas_ro08, u_slot_ro08) are
% exploratory and not in the tables.
% PENDING (owed, 6 runs on today's tree; the lab was taken by another user at ~14:00): re-fly u_meas_t08 (decides the
% falsifier at 8 cm and the measured rim limits: d_tol / d_sep / d_task 8 / 12 / 8, or 6 / 6 / 6); re-fly u_slot_t04
% (the slot rim limits, 4 or 2); u_meas_ri20 x2 and u_slot_ri20 x2 (the pre-registered 20 cm extension radially in).
% Then the reviewer, and U re-flown whole on the final tree.
% Must not say: anything for the rig beyond ~2.5 cm (a real magnet catches only on its 60 mm plate: 16 Oct, three
% slides 3.3-4.3 cm off went past its edge; our contact check confirms within ~24 mm of plate 9); "the model point does
% not matter" along the rim (it matters there, in the slot arm's favour); that the mechanism is shown along the rim; that
% the measured arm's static balance is right at every offset (it is radially in only).

\subsection{The magnet caught off its slot}
\label{sec:res_catch}

This section asks how far off its slot the magnet can catch before the task fails, with the fold-in at the measured
catch point, as the rig folds in, and with the slot assumed instead. It is a twin result. A real magnet catches only on
its 60\,mm plate, and the contact check confirms a catch only within 24\,mm of the plate point
(Table~\ref{tab:constants}); the twin's magnet welds in the air wherever it is told. Up to about 2.5\,cm the sweep is
therefore the rig's case, a real catch error folded in as the rig folds it in. Beyond that it measures the fold-in's
model-error tolerance: how far the attach point the planner is given may lie from the true one before the fold-in
fails, not how far off a real magnet catches. The one real off-slot catch in the sweep is the centre of the next plate
along the rim, 11.6\,cm away: a catch on the wrong plate, which the contact check would refuse on the rig.

\paragraph{How it was flown.} Three carriers on plates 3, 7 and 11 lift the 0.86\,kg ring on 1\,m cables and hold it at
1.0\,m. The newcomer flies this work's own approach to plate 9, which exists only in the twin
(Appendix~\ref{sec:app_proximity_weld}), and the twin's magnet welds its rod's end once it is within 2.5\,cm of the aim
and slower than 0.05\,m/s for 0.15\,s. The aim and the weld point move together by the offset $d$, along the rim
(clockwise, toward the carrier on plate 11) or radially inward, from 0 to 16\,cm. Two arms fold the newcomer in at the
weld: the measured arm at the caught point, as the rig does, and the slot arm at the slot, while the rod pulls from the
caught point, an unmodelled moment that grows with $d$. Both take the caught tip's height, so the arms differ only in
where the planner places the attach point in the ring's plane. The four UAVs then hover, fly one figure-eight
($a = 1.0$\,m) at 0.2\,m/s and land. Two runs per cell, on one host and one code tree.

\paragraph{The bars.} A run fails the task on any of: no weld within 60\,s of the magnet's switch-on, no fold-in within
2\,s of the weld, an abort, a drop, the figure-eight not completed before LAND, or a UAV not landed. The fold-in's bar is
the ring's tilt peak, at most 5$^\circ$ from 1\,s before the weld to 15\,s after it. The arms are compared on
$d_\mathrm{tol}$, the largest offset up to which every cell completed within the bar in both runs. The rule that the
newcomer must stay 0.30\,m from every carrier is scored apart, as $d_\mathrm{sep}$: where the newcomer flies is set by
the arm itself, and the planner has no separation constraint, so crowding is a finding of its own and not a test of
the balance the arms are compared on. The task's limit, $d_\mathrm{task}$, is the smaller of the two. Because equal
limits would show nothing, a mechanism bar asks that the slot arm's excess tilt, at the fold-in or in the steady hover,
grow beyond the repeat spread at $d = 0$ at the two largest offsets both arms flew. All of it was frozen before the
runs.

\begin{table}[htbp]
\centering
\footnotesize
\setlength{\tabcolsep}{3pt}
\caption[The magnet caught off its slot, along the rim]{The magnet caught off its slot along the rim, clockwise toward
the carrier on plate 11, in the rig twin: two runs per cell (a / b). Fold-in: the ring's tilt peak from 1\,s before the
weld to 15\,s after it (bar 5$^\circ$). Steady: the ring's mean tilt in the four-UAV hover, from 15\,s after the weld to
the path. Separation: the newcomer's body to the nearest carrier's, from the weld to LAND (line 0.30\,m). Bold: over the
bar or under the line. Every weld caught within 2.0\,cm of its aim; in every run the tip was within 1\,cm of the aim
19.0--21.1\,s after the magnet's switch-on, and the weld then waited 2.1--48.3\,s for the trigger's speed gate.
$^a$\,Void: a fault of the test runner after the landing (as flown 1.65$^\circ$, 0.62\,m); re-fly owed. $^b$\,No weld
within 60\,s: a task failure. $^c$\,Void: the weld came too late to settle before the path (as flown 2.41$^\circ$,
0.49\,m); re-fly owed. $^d$\,The centre of plate 10: a catch on the wrong plate. $^e$\,This run never settled: the ring
wobbled 3--13$^\circ$ and 6--18\,cm off its reference through the hover and the figure-eight. $^f$\,Welded 68.9\,s after
the switch-on, over the 60\,s allowed: a task failure; its fold-in ran into the path.}
\label{tab:catch_rim}
\begin{tabular}{lcccccc}
\toprule
 & \multicolumn{3}{c}{Measured catch point} & \multicolumn{3}{c}{Slot assumed} \\
\cmidrule(lr){2-4}\cmidrule(lr){5-7}
$d$ (cm) & Fold-in ($^\circ$) & Steady ($^\circ$) & Separation (m) & Fold-in ($^\circ$) & Steady ($^\circ$) & Separation (m) \\
\midrule
0 & 2.21 / 2.02 & 2.29 / 1.93 & 0.65 / 0.65 & 1.58 / 1.70 & 1.02 / 1.46 & 0.65 / 0.65 \\
2 & 1.93 / 1.75 & 2.25 / 1.62 & 0.61 / 0.62 & 1.71 / 1.66 & 1.52 / 1.63 & 0.63 / 0.63 \\
4 & 1.83 / 1.76 & 4.04 / 2.47 & 0.58 / 0.58 & void$^a$ / 1.90 & -- / 2.13 & -- / 0.62 \\
6 & 2.06 / 1.88 & 4.27 / 5.39 & 0.54 / 0.54 & no weld$^b$ / 1.93 & -- / 2.30 & -- / 0.60 \\
8 & void$^c$ / 2.33 & -- / 6.18 & -- / 0.48 & 3.75 / 2.67 & 4.96 / 3.98 & 0.59 / 0.59 \\
11.6$^d$ & \textbf{7.57} / \textbf{7.62} & 5.91 / 6.19 & 0.37 / 0.37 & 3.50 / 3.60 & 4.63 / 5.72 & 0.56 / 0.56 \\
12 & \textbf{5.82} / \textbf{5.14} & 6.36 / 4.93$^e$ & 0.36 / 0.32 & 3.78 / 4.65 & 5.69 / 6.92 & 0.56 / 0.56 \\
16 & \textbf{7.96} / \textbf{8.87} & 7.60 / 5.01 & \textbf{0.22} / \textbf{0.23} & 3.69 / \textbf{7.09}$^f$ & 6.27 / -- & 0.54 / 0.54 \\
\bottomrule
\end{tabular}
\end{table}
% Measured | slot, by d: 0 R2629/R2657 | R2630/R2658; 2 R2631/R2659 | R2632/R2660; 4 R2633/R2661 | R2634/R2662; 6
% R2635/R2663 | R2636/R2664; 8 R2637/R2665 | R2638/R2666; 11.6 R2639/R2667 | R2640/R2668; 12 R2641/R2669 | R2642/R2670;
% 16 R2643/R2671 | R2644/R2672. True offset (cm), measured: 0.7/0.5, 2.3/2.4, 3.8/4.1, 5.9/5.9, 7.8/7.7, 11.5/11.7,
% 12.0/11.9, 15.8/15.7; slot: 1.1/0.1, 2.2/1.7, 4.0/4.0, -/5.5, 8.5/8.1, 11.9/12.0, 12.3/12.6, 15.0/15.9. Model error xy:
% measured 0.1-0.8 cm (the tip's miss of the aim inside the trigger sphere); slot = the true offset. Carrier dip (card
% observation 1, beside the 11.6-16 cm fold-in peaks): slot 36-48 mm (completed runs), measured 85-127 mm. The section 2 static check predicted the
% measured arm's planned body 0.73, 0.67, 0.61, 0.49, 0.36 and 0.23 m from the plate-11 carrier at 0, 2, 4, 8, 12 and 16 cm.

Along the rim the measured arm completed every cell to 6\,cm, and 8\,cm in its one counted run, with its fold-in well
inside the bar (Table~\ref{tab:catch_rim}). From the next plate's centre on, at 11.6, 12 and 16\,cm, its fold-in tilted
the ring past the bar in every run, and at 16\,cm its newcomer came within 0.30\,m of the carrier on plate 11, as the
static check before the runs had predicted. Its steady tilt grew with the offset although its attach point stayed within
a few millimetres of the true one, so the static balance the hypothesis expected at any offset held radially but not
along the rim; why was not separated. The slot arm folds in more gently along the rim: wherever it welded in time its
fold-in stayed under the bar, with a smaller dip of the carriers, and it kept its newcomer near plate 9, clear of the
carriers. Its limit along the rim was set not by the fold-in but by the twin's magnet: at 6\,cm one run never welded,
its tip near the aim but never slow enough for the trigger's speed gate, and at 16\,cm one welded too late and ran into
the path. Those belong to the trigger, which the two arms share until the fold-in, and the rig's magnet has no speed
gate.

\begin{table}[htbp]
\centering
\footnotesize
\setlength{\tabcolsep}{3pt}
\caption[The magnet caught off its slot, radially in]{The magnet caught off its slot radially inward, in the rig twin,
as Table~\ref{tab:catch_rim}. Steady offset: the ring's mean distance from its reference in the four-UAV hover (mm).
Every radial run kept the newcomer at least 0.64\,m from every carrier. The slot arm's newcomer peaked at 0.78--0.79
throttle, against its cap of 0.8, in every radial run with an offset.}
\label{tab:catch_in}
\begin{tabular}{lcccccc}
\toprule
 & \multicolumn{3}{c}{Measured catch point} & \multicolumn{3}{c}{Slot assumed} \\
\cmidrule(lr){2-4}\cmidrule(lr){5-7}
$d$ (cm) & Fold-in ($^\circ$) & Steady ($^\circ$) & Offset (mm) & Fold-in ($^\circ$) & Steady ($^\circ$) & Offset (mm) \\
\midrule
0 & 2.21 / 2.02 & 2.29 / 1.93 & 2 / 1 & 1.58 / 1.70 & 1.02 / 1.46 & 3 / 1 \\
2 & 1.84 / 1.68 & 1.77 / 1.65 & 1 / 1 & 3.23 / 3.29 & 3.82 / 3.98 & 7 / 8 \\
4 & 2.63 / 1.54 & 3.02 / 1.04 & 5 / 2 & 1.83 / 1.81 & 1.31 / 0.99 & 121 / 75 \\
6 & 1.54 / 1.68 & 0.93 / 1.22 & 1 / 1 & 3.45 / 3.27 & 3.15 / 3.05 & 173 / 170 \\
8 & 1.41 / 1.63 & 0.81 / 1.32 & 2 / 1 & \textbf{5.76} / 4.20 & 5.45 / 3.83 & 237 / 192 \\
12 & 1.01 / 1.39 & 0.30 / 1.04 & 3 / 1 & \textbf{10.38} / \textbf{11.71} & 10.44 / 11.91 & 337 / 353 \\
16 & 1.04 / 0.93 & 0.44 / 0.08 & 1 / 3 & \textbf{16.96} / \textbf{14.99} & 16.43 / 15.00 & 396 / 372 \\
\bottomrule
\end{tabular}
\end{table}
% Measured | slot, by d: 0 as along the rim; 2 R2645/R2673 | R2646/R2674; 4 R2647/R2675 | R2648/R2676; 6 R2649/R2677 |
% R2650/R2678; 8 R2651/R2679 | R2652/R2680; 12 R2653/R2681 | R2654/R2682; 16 R2655/R2683 | R2656/R2684. True offset (cm),
% measured: 2.6/2.5, 5.7/3.8, 5.9/6.3, 7.6/8.5, 10.0/12.9, 16.2/14.6; slot: 1.9/2.0, 4.4/3.4, 5.8/5.7, 7.8/6.4,
% 11.4/12.2, 15.4/14.3. Path error after (clock-free figure-8, mm): measured 14.6-17.1; slot 18.4-19.8 at 2 cm, 60.0-90.8
% at 4, 122-124 at 6, 139-168 at 8, 225-238 at 12, 247-265 at 16.

Radially inward the measured arm completed every cell to the largest offset flown, 16\,cm, with its fold-in well inside
the bar and its steady tilt, from 6\,cm on, lower than at $d = 0$ (Table~\ref{tab:catch_in}). The slot arm held the bar
to 6\,cm. From 8\,cm its fold-in tilted the ring past the bar, in one run at 8\,cm and in both beyond, its ring hung
well off its reference in the hover from 4\,cm on, and its newcomer flew near the throttle cap: the unmodelled moment,
as the hypothesis said. Here the attach point matters, and the mechanism bar shows it.

\begin{table}[htbp]
\centering
\footnotesize
\setlength{\tabcolsep}{3pt}
\caption[The catch offset's limits and verdict]{The limits and the verdict of the catch-offset test against its bars,
frozen before the runs. Limits in cm of the nominal offset, with $d_\mathrm{tol}$ on the true offset, run by run, beside
it. In brackets: the value if the void cell's re-fly completes within the bar (and keeps the separation). The
mechanism: the growth of the slot arm's excess tilt over the measured arm's from $d = 0$, against the larger of the two
arms' repeat spreads at $d = 0$ (0.20$^\circ$ at the fold-in, 0.44$^\circ$ steady), at the two largest offsets both arms
flew; the scorer needs two runs a cell, and reads 11.6\,cm where the slot arm's 16\,cm cell has one steady value.}
\label{tab:catch_verdict}
\begin{tabular}{llccccl}
\toprule
Arm & Direction & $d_\mathrm{tol}$ & on the true offset & $d_\mathrm{sep}$ & $d_\mathrm{task}$ & Still to fly \\
\midrule
Measured & along the rim & 6 [8] & 7.7 & 6 [12] & 6 [8] & the 8\,cm re-fly \\
 & radially in & 16 & 16.2 & 16 & 16 & the 20\,cm extension \\
Slot & along the rim & 2 [4] & 5.5 & 2 [4] & 2 [4] & the 4\,cm re-fly \\
 & radially in & 6 & 6.4 & 16 & 6 & the 20\,cm extension \\
\bottomrule
\end{tabular}

\vspace{1ex}
\begin{tabularx}{\textwidth}{>{\raggedright\arraybackslash}p{4.3cm}>{\raggedright\arraybackslash}X>{\raggedright\arraybackslash}p{2.6cm}}
\toprule
Bar & As flown & Verdict \\
\midrule
Confirms: the measured arm's $d_\mathrm{task}$ at least 12\,cm along the rim & 6, at most 8: the fold-in bar from 11.6\,cm
  & not met \\
Confirms: the measured arm's $d_\mathrm{task}$ at least 16\,cm radially in & 16 & met \\
Confirms: the measured arm's $d_\mathrm{tol}$ at least the slot arm's, both directions & along the rim 6 against 2;
  radially 16 against 6 & met; along the rim only through the slot arm's no-weld \\
Confirms: the mechanism shown, both directions & radially: fold-in $+$10.3 and $+$15.5$^\circ$, steady $+$11.4 and
  $+$16.3$^\circ$ at 12 and 16\,cm; along the rim: fold-in $-$0.8 and $-$2.6$^\circ$, steady $+$1.5$^\circ$ at 12\,cm and
  $-$0.01$^\circ$ at 11.6\,cm & radially shown; along the rim not shown \\
Falsifies: the measured arm's $d_\mathrm{task}$ under 8\,cm, either direction & no failure, crowding or tilt over the bar
  at 8\,cm or less (fold-in at most 2.63$^\circ$, separation at least 0.48\,m) & not fired; 8\,cm along the rim has one
  counted run \\
Falsifies: its fold-in over 5$^\circ$ at 4\,cm or less & at most 2.63$^\circ$ & not fired \\
Falsifies: the slot arm's $d_\mathrm{tol}$ over the measured arm's & along the rim 2 against 6; radially 6 against 16
  & not fired \\
\midrule
Overall & & between: neither confirmed nor falsified \\
\bottomrule
\end{tabularx}
\end{table}
% Card Results, "The limits" and "Verdict per frozen bar". Mechanism along the rim by the card's words ("cell mean",
% failed runs included) would read shown (+1.52 at 12 cm, +0.83 at 16 cm on a one-run slot cell); Wesley 23 Oct: the
% scorer's code stands, NOT SHOWN.

Against its bars the result is between: neither confirmed nor falsified (Table~\ref{tab:catch_verdict}). Radially
inward the measured arm reached the 16\,cm asked of it, beat the slot arm, and the mechanism is shown. Along the rim it
fell short of the 12\,cm asked, on the fold-in bar from the next plate's centre on, and the mechanism is not shown: at
the fold-in the slot arm tilted the ring less than the measured arm, and its steady excess did not grow past the repeat
spread at both offsets the scorer read. The measured arm's lead along the rim rests on the slot arm's no-weld at 6\,cm,
which is the trigger's and not the arm's; on the fold-in bar alone the slot arm holds to 12\,cm and the measured arm to
8\,cm. No falsifier fired, though the one at 8\,cm waits for that cell's re-fly.

The twin flatters the attach: its rods are rigid and can push, there is no rotor wash and no contact between rods, and
its magnet welds in the air wherever it is told. At 16\,cm along the rim the newcomer's rod meets the ring about 7\,cm
from the next carrier's, which the rig's rods would not allow.

Six runs are owed and still to fly: a re-fly of each void cell, the measured arm at 8\,cm along the rim, which decides
the falsifier there and the measured arm's limits along the rim, and the slot arm at 4\,cm, which decides the slot
arm's; and the 20\,cm extension radially inward, two runs per arm, pre-registered for a direction in which neither arm
failed by 16\,cm. They fly on a newer code tree than the grid's 54 counted runs, whose flight code differs only in parts
that no run of this test executes; the whole table is re-flown on the final tree before the reviewer sees it.

\pendingfig{fig:catch_offset}{The magnet caught off its slot in the rig twin: the fold-in's ring tilt peak against the
true catch offset, (a) along the rim and (b) radially inward, both arms and both runs of each cell, against the
5$^\circ$ bar, with the 2.5\,cm a real magnet can catch within shaded and the wrong plate's centre marked; (c) along the
rim, the newcomer's closest approach to a carrier against the 0.30\,m line.}{the grid with the two re-flies and the
20\,cm extension}
% Data: tools/catch_offset.py (its JSON for the figure) on R2629-R2684 and the six owed runs. Tool: a new
% thesis_sim_figures function through the thesis-figure skill (Wesley reviews it first).

\paragraph{Answer.} In the twin, with the fold-in at the measured catch point as the rig does it, the task completes
radially inward to the largest offset flown, 16\,cm, and along the rim to 6\,cm, or 8\,cm if that cell's re-fly
completes: from the next plate's centre on, the fold-in tilts the ring past its bar, and at 16\,cm the newcomer crowds
the next carrier. Within the rig's own catch errors, up to about 2.5\,cm, both arms fold in inside the bar. Beyond that
the sweep measures the fold-in's model-error tolerance rather than a real catch: radially the measured point matters,
and with the slot assumed the fold-in fails from 8\,cm; along the rim the mechanism is not shown, and the slot arm folds
in more gently. The result is between, neither confirmed nor falsified, with six runs owed.
